\chapter{Generalidades de la Empresa}
\label{cp:inteligencia-artificial}

\section{Perfil de la Organización}
\label{sec:concepto-ramas-IA}
El termino , a pesar que no posea una definición única y aceptada por todos, la mayor parte de los 
investigadores en este campo concuerdan en un grupo de características fundamentales que distinguen a 
estos sistemas: la capacidad de aprender y adaptarse, solucionar problemas complejos, entender el entorno 
y razonar con eficacia incluso en situaciones inciertas. La  es un ámbito que se alimenta de la 
convergencia de diversas disciplinas de investigación, lo que hace posible que no exista una única 
definición correcta.  Por ende, es crucial que cada investigador establezca de manera precisa su 
perspectiva personal sobre la , pues las diversas interpretaciones conllevan consecuencias 
significativas tanto a nivel práctico como teórico.  En realidad, la selección de un concepto inicial 
determinará el camino y el resultado final de toda la investigación. A pesar de 
sus variaciones, todas las perspectivas acerca de la inteligencia artificial comparten un punto clave: la 
disciplina se enfoca en analizar, diseñar y crear agentes inteligentes que tienen la habilidad de cumplir 
objetivos o metas definidas previamente.

La no es solamente un progreso a nivel técnico, sino que también conlleva una transformación cognitiva 
para las sociedades actuales. No solo se trata de automatizar procesos al adoptar esto, sino que también 
significa una nueva manera de concebir los sistemas complejos como el tránsito urbano bajo un enfoque 
predictivo y adaptativo. Este cambio de paradigma, que va más allá de la programación básica, modifica el 
rol del ser humano como desarrollador de sistemas inteligentes que adquieren conocimiento a partir de la 
experiencia, en lugar de ser solo un operador de aparatos fijos.
La se divide en muchas ramas y subcampos, que se enfocan en diferentes tipos de tareas o metodologías. 
Las principales ramas o subcampos de la son:

\section{Estructura del Departamento de TI}

\section{Infraestructura Tecnológica Base}
